Conformal prediction constructs prediction sets by calibrating a scalar measure of discrepancy between predictions and observations~\cite{shafer2008tutorial}.
This quantity is referred to as a nonconformity score.
Let $\widehat{y}_t$ denote a point prediction of an outcome $y_t$, and define
\begin{equation}
    s_t=S(y_t, \widehat{y}_t),
\end{equation}
where $S$ is an application-dependent score function.
Let $q_t(\beta)$ denote the empirical $(1-\beta)$-quantile of the previously observed scores with the standard finite-sample correction.
The corresponding prediction set is
\begin{equation}
    \mathcal{C}_t(\beta)=\{
    y:S(y, \widehat{y}_t)\leq q_t(\beta) \}.
    \label{eq:conformal_set}
\end{equation}
Under exchangeability of the calibration and test scores, $\mathcal{C}_t(\beta)$ provides marginal coverage of at least $1-\beta$.

In sequential robotic operation, the score distribution may vary over time and exhibit temporal dependence.
Adaptive conformal inference adjusts the effective miscoverage level using observed violations~\cite{gibbs2021adaptive}.
Let $e_t=\mathbbm{1}\{y_t\notin\mathcal{C}_t(\beta_t)\}$ denote a prediction-set violation.
For a target violation level $\delta$, the adaptive update is
\begin{equation}
    \beta_{t+1} = \beta_t + \eta(\delta-e_t),
    \label{eq:adaptive_conformal_update}
\end{equation}
where $\eta>0$ is the adaptation rate.
A violation decreases $\delta_t$ and enlarges subsequent prediction sets, whereas repeated coverage permits tighter sets.
This feedback controls the long-run empirical violation frequency rather than guaranteeing coverage for every individual prediction.

Adaptive conformal methods have been applied to motion planning and safety-critical control to calibrate trajectory-prediction and safety-assessment errors online~\cite{zhou2024safety, huriot2026safe}.
In~this work, we use the upper-tail prediction error of a future energy requirement as the score.
Its calibrated quantile yields a one-sided upper bound used to define the energy-sufficiency certificate and precedence margin.